\section*{Abstract}
Dieser Bericht untersucht die Optimierung quaternionwertiger neuronaler Netzwerke auf vier nichtlinearen vierdimensionalen Regressionsbenchmarks: Ackley, Griewank, Rastrigin und Schwefel. Der Fokus liegt nicht auf einer neuen Modellarchitektur allein, sondern auf der Frage, welche Optimierungsdynamik entsteht, wenn klassische Quaternion-Backpropagation mit Predictive Coding verglichen wird. Als Baselines werden CMLP und GHR reproduziert; danach werden ein reellwertiges Predictive Coding Network (PCN) und eine split-quaternion Variante (CPCN) untersucht. CMLP bleibt auf Ackley, Griewank und Rastrigin die stärkste Baseline, waehrend GHR auf Schwefel leicht besser abschneidet. Das reellwertige PCN reduziert zwar teilweise interne Rekonstruktionsfehler, liefert in der freien Test-Inferenz aber keine brauchbare Funktionsapproximation. CPCN ist numerisch stabiler und erreicht kleinere Fehler als das reellwertige PCN, erklaert auf der Testmenge jedoch noch fast keine Varianz. Das Resultat ist damit folgendes: In $\mathbb{H}$ bzw. in der $\mathbb{R}^4$-Darstellung sind Lernrate, Batchgroesse, Produktreihenfolge, latente Inferenz und Testmodus stark gekoppelt. Predictive Coding ist für QVNNs implementierbar, führt aber in der aktuellen Konfiguration nicht automatisch zu konkurrenzfähiger Generalisierun und benötigt warscheinlich eine grosse Hyperparametersuche.
\vspace{2ex}
